% sup4_09_S14_4: Appendix on Ba3 Prop. 1, part 4, step S14.4: the closure of the classes under the B~'-expansion
% ("which can be decomposed into a sum of renormalized classes"). Input from a08_balaban_prop1.tex.

\subsubsection{S14.4: the classes after the $\tilde B'$-expansion, \cite[p.~433]{Ba3}}\label{sup4:ssec:S144}

\cite[p.~433]{Ba3}: the $\tilde B'$-expansion of a renormalised class gives an expression ``which can be decomposed into a sum of renormalized classes''. Part III proves (S14) that the untouched divergent subgraphs of an expanded graph are renormalised by terms of the expansion of the class, and that a touched divergent subgraph changes type: a $2\phi$ subgraph touched once becomes a $2\phi1A$ subgraph of degree $0$, which needs the partner with its two external scalar legs exchanged; a $1A$ subgraph touched once, or a vacuum component touched twice, becomes a $2A$ subgraph of degree $\le0$, which needs its Ward--Takahashi class. The closure property this requires of the class family is proved here.

\paragraph{Classes.} Localisation labels are attached to each leg of a vector field (the support box of its partition function, a grid box of side $\le2$), to each \cite[(1.6)]{Ba3}/\cite[(1.7)]{Ba3} vertex (its unit cube) and to each block factor \cite[(1.13)--(1.15)]{Ba3} (the block $B^k(y)$); the counterterm graph $T_\gamma G$ of a $2\phi$ subgraph $\gamma$ carries the labels of $G$ (\cite[p.~417]{Ba3}: the legs are attached at one point, the internal vertices keep their localisations). Two graphs are ``in agreement'' \cite[p.~432]{Ba3} if they carry the same label multiset, and the \emph{per-label full class} of a given external type and order is the set of all connected graphs of that type and order with a given label multiset; equivalently, with the partition functions written as multipliers $t_yH_y$ on the vector legs and $t_Q\mathbf 1_Q$ on the \cite[(1.6)]{Ba3}/\cite[(1.7)]{Ba3} vertices, it is the coefficient of one monomial $\prod t$ in the expansion of the generating functional \cite[(1.5)]{Ba3}. This is the class choice of \cref{sup4:lem:F4}, permitted by \cite[p.~431]{Ba3}, and it is the one \cref{app:a06_BALtorus} uses for the regrouping after the cut and for the renormalisation. After the $\tilde B'$-expansion the new vertices are labelled in the same way (one label per $\tilde B'$-leg).

\begin{lemma}[exchange of the external scalar legs]\label{sup4:lem:S1441}
A per-label full class is closed under the exchange of the two external scalar legs of any $2\phi$ or $\phi\phi A$ subgraph, the whole graph included.
\end{lemma}

\begin{proof}
The class is a $t$-monomial coefficient of \cite[(1.5)]{Ba3}, and the Wick expansion sums over all pairings of legs. Exchanging which of two legs is paired with which outside slot is a bijection of pairings; it preserves orders and labels (labels sit on legs and vertices, not on pairings), and it maps graphs of the class to graphs of the class. For the whole graph the two external slots are the arguments of the external field $\Phi_{\mathrm{ext}}(x,x')$, and the sum over pairings contains both assignments.
\end{proof}

\begin{lemma}[the $\tilde B'$-Taylor coefficients of a class]\label{sup4:lem:S1442}
Let $\mathcal C$ be the per-label full class of type $1A$ (resp.\ vacuum) and order $o$, at background $\tilde B=B_0+\tilde B'$ on $\square$. Its first (resp.\ second) $\tilde B'$-Taylor coefficient is the per-label full class of type $2A$ and the same order, with one (resp.\ two) external legs equal to $\tilde B'$ carrying their partition labels; the combinatorial weights are those of the Taylor coefficients.
\end{lemma}

\begin{proof}
On $\square$ the action $S(\phi;B_0+\tilde B'+X)$ depends on $\tilde B'$ and on the other gauge-field legs $X$ ($A'$ with its $g$ and partitions, $\tilde A$) only through the total field $B_0+\tilde B'+X$ (\cite[(1.4)]{Ba3}: both the kinetic term and the averaging term depend on the total field only), and $\tilde B'$-legs carry partition labels like every vector leg. Let $W(\tilde B',X)\defeq\log$ of the Gaussian integral at background $B_0$; $W$ is a function of $\tilde B'+X$, so $\partial/\partial\tilde B'(y)=\partial/\partial X(y)$ on every Taylor coefficient, label by label. By the linked-cluster theorem (\cref{sup2:lem:S8torus}) the Taylor coefficients of $W$ are the sums of all connected graphs, whose vertices are the joint Taylor coefficients of $S$ in the total field and whose lines are $G_k(\square,B_0)$. The $\tilde B'$-expansion of part III produces exactly these joint Taylor coefficients: insertions on lines are the Taylor coefficients of $G_k(\square,B_0+\tilde B')$, and expanded old vertices are the $\tilde B'$-Taylor coefficients of the old vertices; the factorised form of part III and the resolvent expansion \cite[(3.44)]{Ba1} are two groupings of the same polynomials. So the first $\tilde B'$-coefficient of the $1A$ class is the coefficient of one $X$-leg and one $\tilde B'$-leg of $\partial^2W/\partial X^2$, which is the $2A$ class, and the second coefficient of the vacuum class is the same with two $\tilde B'$-legs. (The parity statement that odd total orders in $(A,\tilde B')$ vanish plays no role here: the bijection holds for the full polynomials.)
\end{proof}

\begin{lemma}[compatibility with the cut and with the renormalisation]\label{sup4:lem:S1443}
\textup{(a)} The expanded class is again a per-label full class. Its cut on $\square$ (the lines between different clusters of label boxes become external fields) is the cluster cut of \cref{app:a06_BALtorus} on its label boxes, and the regrouping theorem there applies verbatim. A touched $2\phi$ subgraph and its mirror (\cref{sup4:lem:S1441}) have the same labels, as do a touched $1A$ subgraph and its $2A$ completion (\cref{sup4:lem:S1442}); so partners stay in one block of the cut or are split together. \textup{(b)} The renormalisation step (\cref{sup4:ssec:S21a,sup4:ssec:S21b} and Step~3(D$_4$) of \cref{app:a06_BALtorus}) is applied to per-label full classes, the class choice of \cref{sup4:lem:F4}. The $\tilde B'$-legs are external fields $A$ in the Ward--Takahashi identities (\cref{sup4:lem:S216}), which hold for every external field and every additional field $X$; after the gauge-out S18, $B_0=0$ and the lines are $G_k(0)$.
\end{lemma}

\begin{proof}
(a) The clusters are determined by the label multiset, which the exchange and the completion preserve (the moved ends stay at vertices whose labels lie in the same cluster), so the cut set is unchanged. (b) As stated.
\end{proof}

\paragraph{S14.4, consumed form.} With \cref{sup4:lem:S1441,sup4:lem:S1442,sup4:lem:S1443}, the partners required by the touched subgraphs of part III are terms of the expansion of the class: the exchange partner of a touched $2\phi$ subgraph $S$ is $(S^\pi)_\tau=(S_\tau)^\pi$ with $S^\pi\subset G^\pi$ in the class (\cref{sup4:lem:S1441}), and the $2A$ Ward--Takahashi class at the order of a touched $1A$ subgraph or vacuum component is the set of all single (resp.\ double) $\tilde B'$-attachments to all graphs of its type (\cref{sup4:lem:S1442}). The expansion of a per-label full class is therefore a sum of per-label full classes, which is the decomposition ``into a sum of renormalized classes'' of \cite[p.~433]{Ba3}; S15 (``again a sum of renormalized classes'') inherits it. Ba{\l}aban's one sentence hides a closure property of the class family under differentiation in the background, which the per-label full classes have.
